\subsection{Alcance}

El producto de software a desarrollar es una aplicación de escritorio independiente para Windows 11 orientada a apoyar la toma de decisiones de los comandantes o coordinadores de emergencias durante eventos de incendio en el sector de Aguas Claras, ubicado en los Cerros Orientales de Bogotá.

El sistema integrará una biblioteca estática de escenarios precalculados mediante simulación computacional de dinámica de fluidos (FDS) con una interfaz gráfica georreferenciada. A partir de la ubicación del punto de ignición o borde activo señalado por el usuario y de las condiciones meteorológicas actuales (obtenidas automáticamente vía API externa o ingresadas manualmente), la aplicación consultará y agregará los escenarios compatibles para presentar una estimación visual de la posible propagación del fuego, temperatura, humo y velocidad de expansión.

\textbf{Delimitación del alcance:}
\begin{itemize}
    \item \textbf{Qué hace el producto:} consulta escenarios precalculados almacenados localmente en SQLite, obtiene datos meteorológicos en tiempo real mediante solicitudes HTTP GET a Open-Meteo, agrega resultados compatibles y los despliega sobre un mapa georreferenciado del sector de Aguas Claras.
    \item \textbf{Qué no hace el producto:} no ejecuta simulaciones FDS en tiempo real durante la emergencia, no requiere conexión permanente a internet ni servidores externos propios, no gestiona cuentas de usuario ni roles administrativos, y sus estimaciones constituyen una herramienta de apoyo visual que no sustituye el criterio operativo ni el juicio profesional de los cuerpos de socorro.
\end{itemize}